\section{One Factor Stochastic Interest Rates}
One of the great features of the particle method is that it easily extends to
the case where interest rates are also stochastic. An LSVM with stochastic
interest rates (SIR-LSVM) is defined in a risk-neutral measure  $ \mathbb{Q}  $ by:

\begin{equation}  \label{eq:spot} dS_t=r_tS_t dt + \sigma(t,S_t)a_t S_tdW_t 
\end{equation}

where the short term rate $r_t$ and the stochastic volatility $a_t$ are Ito processes.
\subsection{The Calibration Condition}
In this case the calibration condition reads:

\begin{equation} \label{eq:calibration_condition_spot} \sigma(t,K)^2 \frac{\mathbb{E}^{\mathbb{Q}}[ D_{0t}a_t^2 | S_t = K ]}{ \mathbb{E}^{\mathbb{Q}}[ D_{0t} | S_t = K ] } = \sigma_{Dup}(t,K) ^2 - \frac{ \mathbb{E}^{\mathbb{Q}}[ D_{0t}(r_t-r^0_t)\mathbb{I}_{S_t>K} ]}{ \frac{1}{2}K\partial^2_KC(t,K) },
\end{equation}

where the discount factor $ D_{0t}$ is given by $ D_{0t}=exp\bigg(-\int_0^t r_sds\bigg)$ and

$r_t^0 = \mathbb{E}^{\mathbb{Q}^t} [r_t] = -\partial_t\log P_{0t}$, with $ P_{st} $ the time-$s$ value of the bond of maturity $t$.\\

From the calibration condition (\ref{eq:calibration_condition_spot}), the dynamics of the calibrated SIR-LSVM read as the following nonlinear McKean diffusion: 

$$ dS_t=r_tS_t dt + \sigma(t,S_t,\mathbb{Q}_t)a_t S_tdW_t .$$

Alternatively, it reads as the following nonlinear McKean diffusion for the
forward $f_t = S_t/P_{tT}$ in the forward measure $\mathbb{Q}^T$ , where $T$ denotes the last maturity date for which we want to calibrate the market smile:

\begin{equation}\label{eq: forwardPrice}
\frac{df_t}{f_t}= \sigma(t,P_{tT}f_t,\mathbb{Q}_t^T)a_t dW_t^T -\sigma_P^T(t)dB^T_t
\end{equation} 

where 

\begin{equation}\label{eq: localVolatility_forward}
\sigma(t,K,\mathbb{Q}_t^T)^2=\bigg( \sigma_{Dup}(t,K)^2-P_{0T} \frac{ \mathbb{E}^{\mathbb{Q}^T}[ P_{tT}^{-1}(r_t-r^0_t)\mathbb{I}_{S_t>K} ]}{ \frac{1}{2}K\partial^2_KC(t,K) } \bigg)\cdot  \frac{\mathbb{E}^{\mathbb{Q}^T}[P_{tT}^{-1} | S_t = K ]}{ \mathbb{E}^{\mathbb{Q}^T}[ P_{tT}^{-1}a_t^2 | S_t = K ] }
\end{equation}

$ \sigma_P^T(t) $ is the volatility of the bond $ P_{tT}$, and $B^T_t$ the (possibly multi-dimensional) $\mathbb{Q}^T$-Brownian motion that drives the interest rate curve.

\subsection{The Particle Method}
A \textit{particle} is then described by the vector-valued process $(f_t^i,a^i_t,r_t^i)$, for $1\leq i \leq N$, where $N$ is the number of simulated particles.
Using the representation (\ref{eq: localVolatility_forward}) for the local volatility, we then define:

\begin{equation}\label{eq: discrete_localVol_general}
\sigma_N(t,K)^2= \bigg( \sigma_{Dup}(t,K)^2 - P_{0T}\frac{\frac{1}{N}\sum_{i=1}^N (P_{tT}^{i})^{-1}(r_t^i-r_t^0)\mathbb{I}_{S_t^i>K} }{\frac{1}{2}K\partial^2_KC(t,K)}\bigg)\cdot\frac{\sum_{i=1}^N (P_{tT}^{i})^{-1} \delta_{t,N}(S_t^i-K)}{\sum_{i=1}^N (P_{tT}^{i})^{-1} (a_t^i)^2\delta_{t,N}(S_t^i-K)}
\end{equation}

and simulate:
\begin{equation}\label{eq: discrete_forwardPrice}
df_t^i = f_t^i\sigma_N(t,f_t^i P_{tT}^i)a_t^idW_t^i-f_t^i\sigma_P^{T,i}(t)dB_t^i
\end{equation}

\subsection{Acceleration Techniques}
\begin{itemize}
\item Not necessary to compute the sum in (\ref{eq: discrete_localVol_general}) for all $i$
\item Save time: compute $\sigma_N$ for a grid $G_{K,t}$ of values of $K$ with size $N_{K,t}<N$
\item Then inter- and extra-polate (cubic splines)
\end{itemize}

\smallskip
\begin{itemize}
\item Moreover: a large number of terms in the sum in (\ref{eq: discrete_localVol_general}) is negligible
\item Disregard $f_t^i$. when it is far from $K$, i.e., $\delta_{t,N}(S_t^i-K) < \eta$ ($\eta$ being a threshold)
\item This requires sorting the particles according to the spot value
\item The cost $O(N\ln N)$ is compensated by the acceleration with the grid
\end{itemize}

\subsection{Malliavin representation of the local volatility}\label{subsec: malliavinGeneral}
Thanks to an application of the Clark-Ocone's formula we can give another expression of the contribution of stochastic interest rates to local volatility:
\begin{equation}\label{eq: malliavinRepresentation}
\frac{ \mathbb{E}^{\mathbb{Q}}[ D_{0t}(r_t-r^0_t)\mathbb{I}_{S_t>K} ]}{ \frac{1}{2}K\partial^2_KC(t,K) } = \frac{2}{K}\int _0^t \mathbb{E}^{\mathbb{Q}^t}[\sigma_r^t(s)\cdot \MD _s^{B^t}S_t  |S_t = K ]ds
\end{equation}

where $ \MD_s^{B^t} $ is the Malliavin derivative with respect to the Brownian motion $B^t$.
This brings few different benefits:
\begin{itemize}
\item[1.] instead of an unconditional expectation we get a conditional expectation given $ S_t = K $ (easier for MC);
\item[2.] the term $\partial_K^2C$ cancels out: fortunate case as it is sensitive to the strike interpolation/extrapolation method;
\item[3.] both $\mathbb{E}^{\mathbb{Q}}[ D_{0t}(r_t-r^0_t)\mathbb{I}_{S_t>K} ]$ and $K\partial_K^2C$ are very small for strikes $K$ that are far from the money. Hence we get rid of a problematic $0/0$ ratio.
\end{itemize}
By specifying a particular model for the interest rate and computing the Malliavin derivative in (\ref{eq: malliavinRepresentation}), one gets more explicit formulas for the interest rate contribution to the local volatility.
We shall see this applied to case of the Hull-White model.
